\documentclass[11pt,a4paper]{ctexart}
\usepackage[margin=2.45cm]{geometry}
\usepackage{amsmath,amssymb,amsthm,mathtools}
\usepackage{booktabs}
\usepackage[hidelinks]{hyperref}
\usepackage{microtype}

\newtheorem{theorem}{定理}
\newtheorem{lemma}[theorem]{引理}
\newtheorem{proposition}[theorem]{命题}
\theoremstyle{remark}
\newtheorem{remark}[theorem]{注}
\newcommand{\Q}{\mathbb Q}
\newcommand{\Z}{\mathbb Z}
\newcommand{\Norm}{\operatorname N}
\newcommand{\vp}{v_{\mathfrak p}}

\title{方程 \(2^n-n=k^2\) 的一个显式上界\\
  \large 实对数、\(2\)-进线性形式及 LLL 约化的探索}
\author{}
\date{2026年8月}

\begin{document}
\maketitle

\begin{abstract}
本文把基于 Ridout 定理的非有效有限性证明替换为一个有效论证。
通过在 \(\Z[\sqrt2]\) 中用基本单位 \(3+2\sqrt2\) 约化，并应用
Matveev 的显式对数线性形式下界，证明每个正整数解均满足
\[
 n\le 4\,999\,999\,999\,999\,999.
\]
该常数很粗，但是真正可计算的。文末给出精确的 \(2\)-进线性形式，
说明 Yu 型估计以及 LLL 和 Baker--Davenport 约化应如何介入，同时指出
非齐次项随约化因子变化所造成的主要困难。
\end{abstract}

\section{问题与初等限制}

考虑
\begin{equation}\label{eq:main}
 2^n-n=k^2,\quad n,k\in\Z_{>0}.
\end{equation}

若 \(n\ge4\) 为偶数，令 \(x=2^{n/2}\)，则
\[
 (x-k)(x+k)=n.
\]
两因子均为正整数，故 \(x+k\le n\)；但 \(x\ge n\)，矛盾。
直接检查 \(n=2\) 也不成立。因此除解 \((n,k)=(1,1)\) 外，可写
\begin{equation}\label{eq:mq}
 n=2m+1,\quad m\ge1,\quad q=2^m.
\end{equation}
原方程化为
\begin{equation}\label{eq:norm}
 k^2-2q^2=-n.
\end{equation}

还有两个便于后续筛选的必要条件。首先，奇数平方模 \(8\) 等于 \(1\)，
故 \(n\ge3\) 时
\begin{equation}\label{eq:mod8}
 n\equiv7\pmod 8.
\end{equation}
其次，若奇素数 \(p\mid n\)，则 \(p\nmid k\)，且
\(k^2\equiv2^n\pmod p\)。因为 \(n\) 为奇数，Legendre 符号给出
\[
 1=\left(\frac{2^n}{p}\right)=\left(\frac2p\right),
\]
所以每个素因子均满足
\begin{equation}\label{eq:primefilter}
 p\equiv1\ \text{或}\ 7\pmod8.
\end{equation}

\section{在 \texorpdfstring{\(\Z[\sqrt2]\)}{Z[sqrt(2)]} 中约化}

记
\[
 A=k+q\sqrt2,\quad A'=k-q\sqrt2,\quad
 \eta=3+2\sqrt2.
\]
由 \eqref{eq:norm}，有 \(AA'=-n\)。这里 \(A>0>A'\)，而
\(\eta\) 是 \(\Q(\sqrt2)\) 中范数为 \(1\) 的基本正单位。

选择唯一整数 \(t\ge0\)，使
\begin{equation}\label{eq:reduction}
 \beta=A\eta^{-t},\quad
 \sqrt n\le\beta<\eta\sqrt n.
\end{equation}
因为 \(\eta^{-1}=3-2\sqrt2\in\Z[\sqrt2]\)，故
\(\beta\in\Z[\sqrt2]\)。其共轭满足
\begin{equation}\label{eq:beta-conj}
 \beta'=-\frac n\beta,\quad |\beta'|\le\sqrt n.
\end{equation}

此外 \(t<m\)。事实上，
\[
 \eta^t=\frac A\beta
 <\frac{2^{m+3/2}}{\sqrt{2m+1}}<\eta^m.
\]
最后一个不等式在 \(m=1\) 时可直接检查；在 \(m\ge2\) 时由
\(\eta>4\) 立即得到。因此 Matveev 定理中的系数上界可以取 \(B=m\)。

\section{实对数线性形式}

定义
\begin{equation}\label{eq:Delta}
 \Delta=\eta^t2^{-m}\frac{\beta}{2\sqrt2}-1
       =\frac{A}{2q\sqrt2}-1.
\end{equation}
由于 \(k<q\sqrt2\)，有 \(-1<\Delta<0\)，特别地 \(\Delta\ne0\)。
再由 \(A'=-n/A\) 以及 \(A>q\sqrt2\)，得到
\begin{equation}\label{eq:upperDelta}
 0<|\Delta|=\frac{|A'|}{2q\sqrt2}
 =\frac{n}{2q\sqrt2 A}
 <\frac{n}{4q^2}=\frac{n}{2^{2m+2}}.
\end{equation}

以下使用绝对对数高度 \(h(\cdot)\)。由 \eqref{eq:reduction}--
\eqref{eq:beta-conj} 以及 \(\beta\) 为二次代数整数，
\begin{equation}\label{eq:hbeta}
 h(\beta)
 =\frac12\left(\log\max\{1,\beta\}
       +\log\max\{1,|\beta'|\}\right)
 \le\frac{\log n+\log\eta}{2}.
\end{equation}
令 \(\gamma=\beta/(2\sqrt2)\)。高度的次可加性给出
\begin{equation}\label{eq:hgamma}
 2h(\gamma)
 \le \log n+\log\eta+2\log(2\sqrt2).
\end{equation}
记
\begin{equation}\label{eq:c0}
 c_0=\log\eta+2\log(2\sqrt2)
     =3.842188\ldots<3.843.
\end{equation}

我们采用如下标准版本的 Matveev 显式估计：若 \(D\) 次数域中的正实
代数数 \(\alpha_1,\ldots,\alpha_s\) 与整数 \(b_i\) 使
\(\alpha_1^{b_1}\cdots\alpha_s^{b_s}-1\ne0\)，并且
\[
 A_i\ge\max\{Dh(\alpha_i),|\log\alpha_i|,0.16\},
 \quad B\ge\max\{1,|b_1|,\ldots,|b_s|\},
\]
则
\begin{align}\label{eq:matveev}
 \log|\alpha_1^{b_1}\cdots\alpha_s^{b_s}-1|
 {}>&-1.4\,30^{s+3}s^{4.5}D^2(1+\log D)(1+\log B)
       A_1\cdots A_s.
\end{align}

在 \eqref{eq:Delta} 中取
\[
 s=3,\quad D=2,\quad
 (\alpha_1,\alpha_2,\alpha_3)=(\eta,2,\gamma),\quad
 (b_1,b_2,b_3)=(t,-m,1),
\]
并取
\[
 A_1=\log\eta,\quad A_2=2\log2,\quad
 A_3=\log n+c_0,\quad B=m.
\]
由 \eqref{eq:hgamma} 可知这些选择合法。于是
\begin{equation}\label{eq:lowerDelta}
 \log|\Delta|>
 -F(1+\log m)(\log n+c_0),
\end{equation}
其中
\begin{align*}
 F={}&1.4\,30^6 3^{4.5}\,4(1+\log2)
       (\log\eta)(2\log2)\\
  ={}&2.369743807897\ldots\times10^{12}
   <2.37\times10^{12}.
\end{align*}

结合 \eqref{eq:upperDelta} 和 \eqref{eq:lowerDelta}，每个解必须满足
\begin{equation}\label{eq:keyineq}
 (2m+2)\log2-\log(2m+1)
 <2.37\times10^{12}(1+\log m)
      \bigl(\log(2m+1)+3.843\bigr).
\end{equation}

\section{数值截断}

令 \(M=2.5\times10^{15}\)，并对实数 \(x\ge M\) 定义
\begin{align*}
 G(x)={}&(2x+2)\log2-\log(2x+1)\\
 &-2.37\times10^{12}(1+\log x)
       \bigl(\log(2x+1)+3.843\bigr).
\end{align*}
直接代入（所有舍入均朝有利于下界的方向）得到
\[
 G(M)>1.05\times10^{13}>0.
\]
同时
\begin{align*}
G'(x)={}&2\log2-\frac2{2x+1}\\
&-2.37\times10^{12}\left(
 \frac{\log(2x+1)+3.843}{x}
 +\frac{2(1+\log x)}{2x+1}\right).
\end{align*}
括号内两项在 \(x\ge M\) 上递减，而其在 \(M\) 处乘以
\(2.37\times10^{12}\) 后小于 \(0.073\)。因此
\(G'(x)>1.31>0\)。所以 \(G(x)>0\) 对全部 \(x\ge M\) 成立，
这与 \eqref{eq:keyineq} 矛盾。

\begin{theorem}\label{thm:bound}
方程 \eqref{eq:main} 的每个正整数解均满足
\[
 \boxed{n\le4\,999\,999\,999\,999\,999}.
\]
\end{theorem}
\begin{proof}
解 \(n=1\) 显然满足结论。其余解有 \(n=2m+1\)、\(m\ge1\)。
上述论证说明 \(m<M=2.5\times10^{15}\)。由于 \(m\) 为整数，
\(m\le2\,499\,999\,999\,999\,999\)，结论随即得到。
\end{proof}

\section{精确的 \texorpdfstring{\(2\)-进}{2-adic} 线性形式}

写
\begin{equation}\label{eq:abXY}
 \beta=a+b\sqrt2,\quad
 \eta^t=X_t+Y_t\sqrt2,\quad a,b,X_t,Y_t\in\Z.
\end{equation}
比较 \(A=\eta^t\beta=k+2^m\sqrt2\) 的 \(\sqrt2\) 系数，得到
\begin{equation}\label{eq:coefficient}
 \boxed{2^m=aY_t+bX_t}.
\end{equation}
这里 \(X_t\) 为奇数、\(Y_t\) 为偶数，且
\begin{equation}\label{eq:vY}
 v_2(Y_t)=v_2(t)+1\quad(t\ge1).
\end{equation}
式 \eqref{eq:vY} 可由
\(Y_t=(\eta^t-\eta^{-t})/(2\sqrt2)\) 及二进制 LTE 得到。
又因 \(2^m\) 为偶数，\eqref{eq:coefficient} 模 \(2\) 表明 \(b\) 为偶数；
范数 \(a^2-2b^2=-n\) 则表明 \(a\) 为奇数。

令 \(\mathfrak p=(\sqrt2)\) 为 \(\Q_2(\sqrt2)\) 的极大理想，规范化
\(\vp(\sqrt2)=1\)，于是 \(\vp(2)=2\)。从
\[
 \beta\eta^t-\beta'\eta^{-t}=2^{m+1}\sqrt2
\]
得到
\begin{equation}\label{eq:padic-product}
 \frac{\beta}{\beta'}\eta^{2t}-1
 =\frac{2^{m+1}\sqrt2\,\eta^t}{\beta'}.
\end{equation}
因为 \(n\) 为奇数，\(\beta'\) 是 \(\mathfrak p\)-进单位，故有精确估值
\begin{equation}\label{eq:exactvp}
 \boxed{
 \vp\!\left(\frac{\beta}{\beta'}\eta^{2t}-1\right)=2m+3.}
\end{equation}

在适当定义的 \(2\)-进对数分支上，\eqref{eq:exactvp} 等价于研究
\begin{equation}\label{eq:padic-linear}
 \Lambda_{\mathfrak p}
 =2t\log_{\mathfrak p}\eta
  +\log_{\mathfrak p}\!\left(\frac\beta{\beta'}\right),
 \quad \vp(\Lambda_{\mathfrak p})=2m+3.
\end{equation}
这给出了 Yu 型 \(p\)-进对数线性形式估计的直接入口。由于
\[
 h(\beta/\beta')\le2h(\beta)\le\log n+\log\eta,
\]
显式 Yu 下界预期给出形如
\[
 2m+3\le C_{\mathrm{Yu}}
 (1+\log t)(\log n+C_1)
\]
的完全有效不等式。它与实 Matveev 不等式同阶；未经逐项代入 Yu 定理的
常数，本文不把它宣称为更小的数值上界。

还有一个便于计算的必要条件。由 \eqref{eq:coefficient}、
\eqref{eq:vY}，若 \(m\ge3\)，高阶消去迫使
\begin{equation}\label{eq:valuation-filter}
 \boxed{v_2(b)=v_2(t)+1.}
\end{equation}
否则 \(aY_t\) 与 \(bX_t\) 的二进制估值不同，其和的估值只是两者中的
较小者，不可能等于 \(m\)。例如已知解 \(n=7\) 对应
\(m=3,t=1,\beta=1+2\sqrt2\)，确有 \(v_2(b)=1=v_2(t)+1\)。

约化条件还给出
\[
 |a|<\frac{\eta+1}{2}\sqrt n,\quad
 0<b<\frac{\eta+1}{2\sqrt2}\sqrt n.
\]
结合定理 \ref{thm:bound}，有 \(b<1.71\times10^8<2^{28}\)。所以
\eqref{eq:valuation-filter} 将候选分成至多 \(27\) 个二进制估值层：
\[
 1\le v_2(b)\le27,\quad 0\le v_2(t)\le26.
\]

\section{LLL/Baker--Davenport 约化：可用形式与障碍}

当 \(m\ge3\) 时，\eqref{eq:upperDelta} 给出 \(|\Delta|<1/2\)。
对 \(\log(1+\Delta)\) 使用 \(|\log(1+x)|<2|x|\)，得到
\begin{equation}\label{eq:BDform}
 \left|t\tau-m+\mu_{a,b}\right|
 <\frac{n}{2^{2m+1}\log2},
\end{equation}
其中
\begin{equation}\label{eq:taumu}
 \tau=\frac{\log\eta}{\log2},\quad
 \mu_{a,b}=\frac{\log\bigl((a+b\sqrt2)/(2\sqrt2)\bigr)}{\log2}.
\end{equation}
这正是 Baker--Davenport（或 Dujella--Peth\H{o}）约化引理处理的形状。
对一个固定的 \(\beta=a+b\sqrt2\)，可取 \(\tau\) 的高阶连分数收敛分母
\(Q>6M\)，计算
\[
 \varepsilon=\|Q\mu_{a,b}\|-M\|Q\tau\|.
\]
若 \(\varepsilon>0\)，则标准约化引理会把 \(m\) 从
\(M=2.5\times10^{15}\) 降到一个很小的对数级界。等价地，可以对
\(\log\eta,\log2,\log(\beta/(2\sqrt2))\) 构造三维 LLL 格。

但是，\(\mu_{a,b}\) 并非固定常数：每个潜在解都有不同的约化因子
\(\beta\)，而现有上界下可能的 \((a,b)\) 太多。因此不能把
\eqref{eq:BDform} 直接交给一次经典 Baker--Davenport 计算并得到统一结论。
这也是当前从 ``显式但巨大'' 到 ``可穷举'' 的主要障碍。

较有希望的后续方案是把实、\(2\)-进信息联合起来：
\begin{enumerate}
 \item 由 \eqref{eq:valuation-filter} 写
 \(b=2^s b_0\)、\(t=2^{s-1}t_0\)，其中 \(b_0,t_0\) 为奇数，
 且 \(1\le s\le27\)。
 \item 在每一层中展开 \eqref{eq:padic-linear}：
 \[
 \log_{\mathfrak p}\frac{a+b\sqrt2}{a-b\sqrt2}
 =2\left(x+\frac{x^3}{3}+\frac{x^5}{5}+\cdots\right),
 \quad x=\frac{b\sqrt2}{a}.
 \]
 由于 \(a\) 奇、\(b\) 偶，该级数具有足够高的 \(\mathfrak p\)-进精度。
 式 \eqref{eq:exactvp} 因而把 \(t_0\) 限制在一个模巨大二次幂的剩余类中。
 \item 同时使用
 \[
 a^2-2b^2=-(2m+1),
 \]
 将 \(m\) 消去，并对变量 \((a,b,t)\) 建立二维或三维的
 \(2\)-进格。先用 \(2\)-进 LLL 排除绝大多数估值层和剩余类，再对
 幸存的有限个 \(\beta\) 使用 \eqref{eq:BDform} 做实 LLL/连分数约化。
\end{enumerate}
该方案已经给出可实施的计算框架，但完成统一的格构造及严格舍入验证
仍需进一步工作，故本文不声称已求出全部解。

\section{结论}

Ridout 方法只说明解集有限；单位约化加显式 Matveev 定理则给出了
可计算的数值上界。当前严格结论是定理 \ref{thm:bound}。精确等式
\eqref{eq:exactvp} 和估值过滤 \eqref{eq:valuation-filter} 表明，
\(2\)-进方法确实含有额外信息；但经典 Baker--Davenport 的非齐次项
随 \(\beta\) 变化，必须先借助 \(2\)-进分层或统一格构造把候选
\(\beta\) 压缩到有限且实际可处理的集合。

\begin{thebibliography}{9}
\bibitem{Matveev}
E.~M. Matveev,
An explicit lower bound for a homogeneous rational linear form in logarithms
of algebraic numbers II,
\emph{Izvestiya: Mathematics} \textbf{64} (2000), 1217--1269.

\bibitem{Yu}
K. Yu,
Linear forms in \(p\)-adic logarithms II,
\emph{Compositio Mathematica} \textbf{74} (1990), 15--113.

\bibitem{BD}
A. Baker and H. Davenport,
The equations \(3x^2-2=y^2\) and \(8x^2-7=z^2\),
\emph{Quarterly Journal of Mathematics} \textbf{20} (1969), 129--137.

\bibitem{LLL}
A. K. Lenstra, H. W. Lenstra Jr., and L. Lov\'asz,
Factoring polynomials with rational coefficients,
\emph{Mathematische Annalen} \textbf{261} (1982), 515--534.
\end{thebibliography}

\end{document}
